%%=============================================================================
%% Methodologie
%%=============================================================================

\chapter{\IfLanguageName{dutch}{Methodologie}{Methodology}}%
\label{ch:methodologie}

\section{\IfLanguageName{dutch}{Onderzoeksopzet}{Research Design}}
Dit onderzoek volgt een experimenteel onderzoeksdesign waarin drie testmomenten met elkaar worden vergeleken: een vroege usability‑test, een late usability‑test en een gecombineerde analyse van beide. Het doel is om vast te stellen of het moment waarop usability‑testing wordt uitgevoerd een significante impact heeft op de gebruiksvriendelijkheid van een React‑gebaseerde webapplicatie. De studie combineert traditionele usability‑metingen met datagedreven telemetry‑analyse via Microsoft Application Insights. Deze gemengde methode maakt het mogelijk om zowel kwalitatieve als kwantitatieve inzichten te verkrijgen en om verschillen tussen testmomenten statistisch te toetsen.
Het onderzoek bestaat uit twee grote componenten. Het technisch component omvat de ontwikkeling van een proof‑of‑concept React‑applicatie met geïntegreerde Application Insights. Het onderzoekend component omvat de uitvoering van de usability‑tests, de dataverzameling, de statistische analyse en de interpretatie van de resultaten. Door deze twee componenten te combineren, wordt een volledig beeld gevormd van zowel het gebruikersgedrag als de impact van ontwerpwijzigingen doorheen het ontwikkelingsproces.


\section{\IfLanguageName{dutch}{Steekproef en rekrutering}{Sampling and Recruitment}}
De steekproef bestaat uit drie onafhankelijke groepen gebruikers: één voor de vroege usability‑test en één voor de late usability‑test en één voor de gecombineerde-test. Elke groep omvat tien tot vijftien deelnemers, wat in lijn ligt met de aanbevelingen van Rubin en Chisnell, die aangeven dat kleine groepen al voldoende zijn om de meeste usability‑problemen te identificeren \autocite{RubinChisnell2008}. De deelnemers worden gerekruteerd via een online oproep op sociale media, waarbij geïnteresseerden zich kunnen registreren via een formulier.
De toewijzing aan de vroege, late of gecombineerde testgroep gebeurt volledig willekeurig. Randomisatie vermindert het risico op systematische verschillen tussen de groepen en verhoogt de interne validiteit van het onderzoek. Na de dataverzameling wordt een Durbin–Watson‑test uitgevoerd om te controleren of de residuen van beide groepen onafhankelijk zijn, wat een belangrijke voorwaarde is voor de ANOVA‑analyse.
Als inclusiecriterium geldt dat deelnemers voldoende Nederlands moeten beheersen. De applicatie en de instructies zijn Nederlandstalig, waardoor taalvaardigheid een potentiële confounder vormt. Door anderstalige deelnemers uit te sluiten, wordt taalgerelateerde bias vermeden en wordt de betrouwbaarheid van de resultaten verhoogd.


\section{\IfLanguageName{dutch}{Technisch component: Proof-of-Concept}{Technical Component: Proof of Concept}}
Het technisch component bestaat uit de ontwikkeling van een React‑gebaseerde webapplicatie met een meerstapsformulier. De applicatie is ontworpen om typische usability‑problemen te simuleren die vaak voorkomen in real‑world webapplicaties, zoals validatiefouten, onduidelijke navigatie en variabele taakcomplexiteit.
Microsoft Application Insights wordt geïntegreerd om gebruikersinteracties automatisch te registreren. De tool verzamelt gegevens over page views, click events, validatiefouten, time‑per‑step, funnels en navigatiepaden \autocite{MicrosoftAI}. Deze data worden gebruikt om objectieve verschillen tussen de testmomenten te analyseren. De configuratie omvat het definiëren van custom events, het opzetten van funnels voor het meerstapsformulier en het exporteren van ruwe data voor statistische analyse.


\section{\IfLanguageName{dutch}{Onderzoekend component: Usability-tests}{Research Component: Usability Testing}}
De usability‑tests worden uitgevoerd in twee fasen. In de vroege testfase doorlopen deelnemers een eerste versie van het meerstapsformulier. Tijdens deze sessies worden taakvoltooiingstijd, fouten, drop‑off‑momenten en beoordelingsscores (in-app schaal 1–5 en NPS) geregistreerd. Daarnaast worden alle interacties automatisch gelogd via Application Insights. De resultaten van deze fase worden gebruikt om verbeteringen aan te brengen in de applicatie.
In de late testfase doorlopen deze deelnemers de verbeterde versie van de applicatie. Dezelfde metingen worden opnieuw uitgevoerd, zodat verschillen tussen de vroege en late testmomenten kunnen worden vastgesteld. 
De gecombineerde analyse omvat een vergelijking van beide testmomenten op basis van zowel traditionele usability‑metingen als telemetry‑data. Deze aanpak maakt het mogelijk om te bepalen welke verbeteringen de grootste impact hebben gehad en hoe gebruikersgedrag evolueert doorheen het ontwikkelingsproces.


\section{\IfLanguageName{dutch}{Onderzoekstechnieken en dataverzameling}{Research Techniques and Data Collection}}
De dataverzameling bestaat uit drie categorieën: traditionele usability‑metingen, telemetry‑data en subjectieve tevredenheidsscores. Traditionele metingen omvatten taakvoltooiingstijd, foutfrequentie en drop‑off‑momenten, zoals beschreven door Sauro en Lewis \autocite{SauroLewis2016}. Telemetry‑data worden verzameld via Application Insights en omvatten page views, click events, validatiefouten en time‑per‑step. De subjectieve tevredenheid van gebruikers wordt gemeten via een in-app beoordelingsschaal (1–5) en de Net Promoter Score (0–10), geïntegreerd in het feedbackscherm van de applicatie. Omwille van deze in-app integratie werd gekozen voor een vereenvoudigde schaal in plaats van de volledige System Usability Scale, wat in lijn is met onderzoek dat kortere schalen als valide alternatief beschouwt bij kleine steekproeven \autocite{SauroLewis2016}.
Alle data worden geanonimiseerd opgeslagen en verwerkt. De ruwe telemetry‑data worden geëxporteerd naar een analyseomgeving waarin statistische tests worden uitgevoerd. De combinatie van kwalitatieve observaties, kwantitatieve metingen en geautomatiseerde telemetry‑data zorgt voor een robuust en multidimensionaal beeld van de gebruiksvriendelijkheid van de applicatie.

\section{\IfLanguageName{dutch}{Statistische analyse}{Statistical Analysis}}
Om te bepalen of er significante verschillen bestaan tussen de vroege en late testmomenten, wordt een eenweg‑ANOVA uitgevoerd. Deze test is geschikt voor het vergelijken van gemiddelden tussen twee of meer onafhankelijke groepen. De nulhypothese stelt dat er geen significant verschil bestaat tussen de testmomenten, terwijl de alternatieve hypothese stelt dat er wel een verschil is. De drempelwaarde voor significantie wordt vastgesteld op $p < 0{,}05$, zoals gebruikelijk in sociaal‑wetenschappelijk onderzoek.
Voorafgaand aan de ANOVA worden de assumpties van normaliteit (Shapiro-Wilk-test), homogeniteit van varianties (Levene-test) en onafhankelijkheid van residuen (Durbin-Watson-test) gecontroleerd. Wanneer de normaliteitsassumptie geschonden is - wat bij Likert-schalen en teldata bij kleine steekproeven doorgaans het geval is \autocite{SauroLewis2016} - wordt de niet-parametrische Kruskal-Wallis-toets als primaire toets gehanteerd. Bij een significant resultaat worden paarsgewijze Mann-Whitney~U-toetsen uitgevoerd met Bonferroni-correctie om te bepalen tussen welke specifieke groepen de verschillen zich bevinden.
Voor categorische uitkomstvariabelen, zoals de voltooiingsratio (formulier wel of niet ingediend), is een ANOVA niet geschikt omdat het om een binaire uitkomst gaat. In dat geval wordt een chi-kwadraattoets ($\chi^2$) uitgevoerd, die nagaat of de verdeling van voltooide en niet-voltooide sessies significant verschilt tussen de drie testgroepen \autocite{Field2018}. De verwachte celfrequenties worden vooraf gecontroleerd; indien een of meer cellen een verwachte frequentie onder vijf hebben, wordt de Fisher-exacttest als alternatief toegepast.
Naast de significantietoets wordt per ANOVA‑analyse het determinatiecoëfficiënt $R^2$ berekend, equivalent aan $\eta^2$ in een eenweg‑ANOVA. $R^2$ geeft aan welk aandeel van de totale variantie in de afhankelijke variabele verklaard wordt door het testmoment als groepsfactor. Volgens de richtlijnen van \textcite{Cohen1988} geldt $R^2 \geq 0{,}01$ als klein effect, $R^2 \geq 0{,}06$ als middelmatig effect en $R^2 \geq 0{,}14$ als groot effect. Door zowel de $p$‑waarde als $R^2$ te rapporteren, wordt niet alleen uitspraak gedaan over de statistische significantie maar ook over de praktische relevantie van de vastgestelde verschillen.

\section{\IfLanguageName{dutch}{Validiteit en betrouwbaarheid}{Validity and Reliability}}
Om de interne validiteit te waarborgen, worden deelnemers willekeurig toegewezen aan de vroege, late of gecombineerde testgroep. De externe validiteit wordt ondersteund door het gebruik van realistische taken en een applicatie die representatief is voor moderne webapplicaties.
Betrouwbaarheid wordt verhoogd door gestandaardiseerde procedures te gebruiken voor alle testsessies. De instructies, taken en meetmethoden zijn identiek voor alle groepen. Telemetry‑data worden automatisch verzameld, waardoor meetfouten worden geminimaliseerd. Mogelijke bias, zoals taalbias, wordt gecontroleerd door enkel Nederlandstalige deelnemers toe te laten.

\subsection{\IfLanguageName{dutch}{Equivalentie van de formulieren}{Form Equivalence}}

Omdat de drie testgroepen elk een afzonderlijk meerstapsformulier gebruiken (form A: reisaanvraag, form B: evenementregistratie, form C: verblijfsaanvraag), is nagegaan of deze formulieren voldoende vergelijkbaar zijn om onderlinge vergelijking van usability-resultaten te rechtvaardigen. Hiervoor werd een objectieve inhoudsanalyse uitgevoerd op basis van structurele kenmerken van de broncode. Tabel~\ref{tab:form-equivalentie} geeft een overzicht van de gemeten kenmerken per formulier.

\begin{table}[h]
  \centering
  \caption[Structurele vergelijking van de drie formulieren]{Objectieve inhoudsanalyse van de drie meerstapsformulieren}
  \label{tab:form-equivalentie}
  \begin{tabular}{lccc}
    \hline
    \textbf{Kenmerk} & \textbf{Form A} & \textbf{Form B} & \textbf{Form C} \\
    \hline
    Onderwerp                        & Reisaanvraag     & Evenement\-registratie & Verblijfsaanvraag \\
    Aantal stappen                   & 9                & 10               & 10               \\
    Totaal formuliervelden           & 55               & 63               & 65               \\
    Conditionele velden              & 17               & 20               & 20               \\
    Mismatch-waarschuwingen          & 8                & 6                & 7                \\
    Gebruikte veldtypes              & 9                & 9                & 9                \\
    Feedbackstappen                  & 3                & 3                & 3                \\
    \hline
  \end{tabular}
\end{table}

Uit de analyse blijkt dat de drie formulieren kwalitatief sterk overeenkomen: alle drie gebruiken dezelfde onderliggende architectuur (FormShell‑component), dezelfde negen veldtypes, identieke feedbacksecties en vergelijkbare patronen van conditionele logica en mismatch-waarschuwingen. De formulieren zijn ook kwantitatief voldoende vergelijkbaar: form A telt 55 invulvelden, form B 63 en form C 65. De onderlinge verschillen in absolute aantallen zijn beperkt en rechtvaardigen geen structurele bedreiging voor de vergelijkbaarheid van de resultaten.

\section{\IfLanguageName{dutch}{Ethische overwegingen}{Ethical Considerations}}
Alle deelnemers worden vooraf geïnformeerd over het doel van het onderzoek, de aard van de taken en de manier waarop hun gegevens worden verwerkt. Deelname is vrijwillig en deelnemers kunnen op elk moment stoppen. Alle data worden geanonimiseerd opgeslagen en uitsluitend gebruikt voor onderzoeksdoeleinden.



